\section{Conclusion}
\label{sec:conclusion}

This paper presents a cleaned-up deterministic substrate theory fragment built around a three-dimensional brane embedded in $\R^4$.
Its main achievement is not a claim of final unification, but a clarification of structure:
a continuum hyperelastic formulation, a cubic-lattice microscopic interpretation, a geometric contraction channel tied to embedding, a Berry/Wilczek--Zee description of transported narrowband sectors, and a localized-mode program stated as a genuine nonlinear bound-state problem.

The most important strategic conclusion is that the next decisive step is computational.
Within the present cubic-lattice hypothesis, the best first target is a proton- or neutron-like three-channel localized state rather than a purely electron-like excitation.
That target directly tests the axis-aligned triplet sector, the role of prestretch in branch mixing, the existence of non-Abelian holonomy diagnostics, and the possibility of fm-scale confinement by substrate self-guidance.

Accordingly, the paper should be read as a sharpened foundation for numerical work.
If the proposed substrate can produce stable baryon-scale triplet states with the expected order of size, energy, and long-range geometric diagnostics, then the broader interpretive program gains plausibility.
If it cannot, then the present version of the substrate model will have been decisively narrowed or ruled out in its current form.
